\documentclass[
  12pt,
  a4paper,
  brazil
]{article}

\usepackage[utf8]{inputenc}
\usepackage[T1]{fontenc}
\usepackage{times}
\usepackage[brazil]{babel}

\usepackage[
  top=3cm, left=3cm, right=2cm, bottom=2cm
]{geometry}

\usepackage{setspace}
\onehalfspacing

\usepackage{indentfirst}
\setlength{\parindent}{1.25cm}

\usepackage{fancyhdr}
\pagestyle{fancy}
\fancyhf{}
\fancyhead[LE,RO]{\thepage}
\fancyhead[RE,LO]{\textit{O Protocolo que Nao Protege}}
\renewcommand{\headrulewidth}{0.5pt}

\usepackage[colorlinks=true, linkcolor=black, citecolor=black, urlcolor=black]{hyperref}
\usepackage{array, booktabs, longtable}
\usepackage{graphicx}
\usepackage[bottom]{footmisc}

\usepackage{titlesec}
\titleformat{\section}{\normalfont\bfseries\large}{\thesection}{0.5em}{}
\titleformat{\subsection}{\normalfont\bfseries\normalsize}{\thesubsection}{0.5em}{}
\titleformat{\subsubsection}{\normalfont\bfseries\normalsize}{\thesubsubsection}{0.5em}{}

\usepackage{tocloft}
\renewcommand{\cftsecleader}{\cftdotfill{\cftdotsep}}

\begin{document}

\begin{titlepage}
  \begin{center}
    \vspace*{3cm}
    {\large\bfseries Prof. Marcelo Claro \& Grupo de Pesquisa\par}
    \vspace{3cm}
    {\LARGE\bfseries O Protocolo que Nao Protege\par}
    \vspace{0.5cm}
    {\large\bfseries Critica Epistemica, Juridica e Comparada a Portaria CNPq n. 2.664/2026\par}
    \vspace{1.5cm}
    {\normalsize Ensaio cientifico apresentado ao Conselho Nacional de Desenvolvimento Cientifico e Tecnologico (CNPq) como contribuicao ao debate sobre regulacao de inteligencia artificial generativa na pesquisa academica.\par}
    \vfill
    {\large 2026\par}
  \end{center}
\end{titlepage}

\tableofcontents
\newpage

\section{Resumo}\label{resumo}

Este ensaio submete a Portaria CNPq nº 2.664/2026 --- que exige
declaração obrigatória de uso de IAG --- a crítica rigorosa em 6
dimensões: técnica, hermenêutica, ética, jurídico-brasileira, comparada
(China, EUA, UE) e substantiva. Demonstra-se que: (a) detectores de IA
têm taxas de falso-positivo inaceitáveis (Weber-Wulff et al., 2023; Liu
et al., 2024); (b) a proporção uso/declaração chega a 40:1 (Liang et
al., 2025); (c) o art. 9º contém 7 ambiguidades que geram insegurança
jurídica; (d) o protocolo conflita com a LGPD, o art. 218 CF e o
princípio da legalidade; (e) diferentemente da China (padrão nacional
obrigatório GB 45438-2025), da UE (AI Act Art. 50) e dos EUA (NSF
TRUST), o modelo brasileiro depende exclusivamente de autodeclaração sem
enforcement técnico. Conclui-se que, embora a Portaria represente avanço
simbólico ao reconhecer a IAG, suas fragilidades superam seus acertos,
operando como instrumento de retardamento burocrático que penaliza
inovadores sem constranger mal-intencionados.

\textbf{Palavras-chave:} Inteligência artificial generativa. Portaria
CNPq nº 2. 664/2026. Integridade científica. Detecção de IA. Direito
administrativo. Regulação comparada. Qualis CAPES. LGPD.

\section{Abstract}\label{abstract}

This essay subjects CNPq Ordinance No.~2.664/2026 --- which mandates
compulsory declaration of Generative AI (GAI) use --- to rigorous
critique across six dimensions: technical, hermeneutic, ethical,
Brazilian legal, comparative (China, USA, EU), and substantive. It
demonstrates that: (a) AI detectors have unacceptable false-positive
rates (Weber-Wulff et al., 2023; Liu et al., 2024); (b) the
use-to-disclosure ratio reaches 40:1 (Liang et al., 2025); (c) Art. 9
contains seven ambiguities generating legal uncertainty; (d) the
protocol conflicts with Brazil's LGPD, Art. 218 of the Federal
Constitution, and the principle of legality; (e) unlike China (mandatory
national standard GB 45438-2025), the EU (AI Act Art. 50), and the USA
(NSF TRUST), the Brazilian model relies exclusively on self-declaration
without technical enforcement.

\textbf{Keywords:} Generative artificial intelligence. CNPq Ordinance
No.~2.664/2026. Scientific integrity. AI detection. Administrative law.
Comparative regulation. Qualis CAPES. LGPD.

\subsection{1. Introdução: O Ano da Virada
Regulatória}\label{introduuxe7uxe3o-o-ano-da-virada-regulatuxf3ria}

O ano de 2026 será lembrado como o momento em que três das maiores
potências científicas do mundo --- Brasil, China e União Europeia ---
convergiram para regular o uso de inteligência artificial generativa
(IAG) na pesquisa acadêmica, cada qual com seu modelo normativo. Os
Estados Unidos, por sua vez, optaram por via distinta, centrada em
segurança de pesquisa e combate à influência estrangeira. Este ensaio
situa a Portaria CNPq nº 2.664/2026 neste cenário comparado para revelar
suas fragilidades e limites.

\textbf{O marco brasileiro.} A Portaria CNPq nº 2.664, de 6 de março de
2026, institui a Política de Integridade na Atividade
Científica{[}\^{}1{]}. Seu art. 9º estabelece três diretrizes centrais:
(i) alínea ``c'': obrigatoriedade de declarar o uso de IAG ``em qualquer
fase do desenvolvimento da pesquisa (concepção, redação, análise de
dados, submissão) especificando {[}\ldots{]} a ferramenta utilizada e a
finalidade''; (ii) alínea ``d'': vedação de ``submissão de conteúdo
gerado por IAG como se fosse de autoria humana''; (iii) alínea ``e'':
vedação ao uso de IAG para elaboração de pareceres
científicos\footnote{As alíneas ``c'', ``d'' e ``e'' do art. 9º
  representam avanço simbólico ao trazer o tema para o campo normativo.
  Contudo, a redação da alínea ``c'' é excessivamente abrangente
  (``qualquer espécie e em qualquer fase''), capturando desde corretores
  ortográficos baseados em IA até modelos generativos, sem distinguir
  níveis de risco --- ao contrário da abordagem baseada em risco do EU
  AI Act.}. Simultaneamente, a CAPES disponibilizou o documento ``A
inteligência artificial na pesquisa e no fomento: desafios e
oportunidades'' (Brasil, 2026)\footnote{BRASIL, A. L. \emph{A
  inteligência artificial na pesquisa e no fomento: desafios e
  oportunidades}. Brasília: CAPES, 2026. Documento interno de circulação
  restrita. Embora relevante por inaugurar o debate institucional na
  CAPES, não tem força normativa e não soluciona o problema central da
  detecção. Disponível mediante solicitação à CAPES. Sua principal
  contribuição é sinalizar que a agência reconhece a IA como tema
  prioritário.}, que reconhece a urgência do debate mas não oferece
solução operacional.

A pergunta fundamental --- ``como provar que alguém usou IA?'' --- não
encontra resposta no texto da Portaria nem na literatura técnica. Este
ensaio sustenta que, ao exigir o que é logicamente impossível de
verificar, o protocolo CNPq opera como instrumento de retardamento
burocrático da pesquisa, penalizando inovadores e não constrangendo
mal-intencionados. {[}\^{}1{]}: BRASIL. Portaria CNPq nº 2.664, de 6 de
março de 2026. Institui a Política de Integridade na Atividade
Científica. \emph{Diário Oficial da União}, Brasília, 11 mar. 2026.
Pontos positivos: reconhece oficialmente a IAG como ferramenta legítima,
não a proíbe, alinhando-se ao movimento internacional. Fragilidade
central: não define sanções específicas nem mecanismos de verificação,
criando norma sem enforcement --- o que o direito administrativo chama
de ``lex imperfecta''. A ausência de tipificação sancionatória viola o
princípio da legalidade estrita (art. 5º, II, CF) ao expor o pesquisador
ao arbítrio interpretativo do agente público.

\subsection{2. Dimensão Técnica: O Problema Insolúvel da
Detecção}\label{dimensuxe3o-tuxe9cnica-o-problema-insoluxfavel-da-detecuxe7uxe3o}

\subsubsection{2.1. A falência dos detectores
automatizados}\label{a-faluxeancia-dos-detectores-automatizados}

O art. 9º, alínea ``c'', pressupõe que o uso não declarado de IAG pode
ser identificado e punido. Esta premissa é falseada por um corpo robusto
de evidências empíricas.

Weber-Wulff et al.~(2023) testaram 14 detectores comerciais de texto
gerado por IA e encontraram ``altas taxas de falso-positivo'',
concluindo que ``nenhum detector é confiável o suficiente para uso em
contextos de alta penalidade''{[}\^{}4{]}. A amostra incluiu textos
acadêmicos reais e gerados por diferentes modelos (GPT-3, GPT-3.5,
GPT-4), em múltiplos idiomas, e os resultados foram consistentes:
nenhuma ferramenta atingiu simultaneamente sensibilidade
\textgreater80\% e especificidade \textgreater90\%.

Liu et al.~(2024) compararam detectores automatizados com julgamento
humano na identificação de textos médicos gerados por LLMs. O estudo,
publicado no \emph{International Journal for Educational Integrity},
concluiu que ``nenhum método atinge confiabilidade aceitável'' e que ``a
concordância entre detectores é baixa''\footnote{LIU, J. Q. J. et
  al.~The great detectives: humans versus AI detectors in catching large
  language model-generated medical writing. \emph{International Journal
  for Educational Integrity}, v. 20, n.~1, 2024. DOI:
  10.1007/s40979-024-00155-6. Originalidade: compara detecção humana
  vs.~automatizada em contexto médico. Limitação: amostra pequena (n=200
  textos). Relevância para o CNPq: demonstra que mesmo especialistas
  humanos não identificam IA de forma confiável, invalidando a premissa
  de que revisores por pares podem fazê-lo.}.

Pudasaini et al.~(2025), em \emph{survey} abrangente no \emph{Journal of
Academic Ethics} (Qualis A1), documentaram que a paráfrase --- técnica
trivial disponível em qualquer LLM --- reduz a acurácia dos detectores a
níveis próximos do aleatório. Os autores testaram 5 detectores contra
textos submetidos a paráfrase, tradução e \emph{prompt injection},
constatando que ``a taxa de detecção cai de 85\% para 32\% após
paráfrase mínima''\footnote{PUDASAINI, S. et al.~Survey on AI-generated
  plagiarism detection: the impact of large language models on academic
  integrity. \emph{Journal of Academic Ethics}, v. 23, p.~1137-1170,
  2025. DOI: 10.1007/s10805-024-09576-x. Qualis A1. Referência
  essencial. Contribuição: mapeia sistematicamente 47 detectores, 23
  datasets e 18 estratégias de evasão. A seção 4.3 sobre ``evasion
  strategies'' deveria ser leitura obrigatória na CAPES. Limitação: foco
  em inglês; idiomas como português têm desempenho de detecção ainda
  pior.}.

O estudo de larga escala mais recente é o de Liang et al.~(2025), que
analisou 164.579 artigos completos e 5,2 milhões de resumos. A conclusão
é devastadora para o protocolo CNPq: entre os 75.172 artigos publicados
desde 2023, apenas 76 (\textasciitilde0,1\%) declararam uso de IA.
Contudo, a proporção estimada de artigos com conteúdo gerado por IA no
primeiro trimestre de 2025 era 40 vezes maior que a taxa de declaração
(40:1)\footnote{LIANG, W. et al.~Discrepancy between AI content
  proportion and disclosure rate in scientific publications.
  \emph{Nature Human Behaviour}, 2025. A escala do estudo (5,2 milhões
  de artigos) lhe confere poder estatístico ímpar. DOI:
  10.48550/arXiv.2512.06705. A razão 40:1 é o dado mais contundente
  contra a eficácia do protocolo brasileiro. Limitação: usa método de
  excesso de vocabulário (Kobak et al.), que fornece estimativas
  probabilísticas em nível de grupo, não diagnósticos individuais;
  possível viés de idioma (inglês super-representado).}. {[}\^{}4{]}:
WEBER-WULFF, D. et al.~Testing of detection tools for AI-generated text.
\emph{International Journal for Educational Integrity}, v. 19, n.~1,
2023. DOI: 10.1007/s40979-023-00146-z. Estudo financiado pelo
Bundesministerium für Bildung und Forschung (BMBF) alémão. Limitação:
testou apenas detectores disponíveis em 2023; versões posteriores
(GPTZero 2.0, Turnitin AI) podem ter desempenho diferente, mas a
literatura de 2024-2025 (Liu et al.; Pudasaini et al.) confirma que o
problema persiste. Forte: metodologia multiplataforma robusta com 14
ferramentas. Fraco: não testou cenário de paráfrase adversarial.

\subsubsection{2.2. A anomalia do art. 9º, alínea
``c''}\label{a-anomalia-do-art.-9uxba-aluxednea-c}

A redação da alínea ``c'' --- ``IAG, de qualquer espécie e em qualquer
fase'' --- padece de \textbf{hipergeneralização}. A IA não é monolítica.
Kocak et al.~(2025) distinguem ao menos quatro níveis de uso de IA na
pesquisa acadêmica, cada qual com diferentes implicações éticas e de
detecção{[}\^{}8{]}:

{\def\LTcaptype{none} % do not increment counter
\begin{longtable}[]{@{}
  >{\raggedright\arraybackslash}p{(\linewidth - 6\tabcolsep) * \real{0.1489}}
  >{\raggedright\arraybackslash}p{(\linewidth - 6\tabcolsep) * \real{0.1915}}
  >{\raggedright\arraybackslash}p{(\linewidth - 6\tabcolsep) * \real{0.2766}}
  >{\raggedright\arraybackslash}p{(\linewidth - 6\tabcolsep) * \real{0.3830}}@{}}
\toprule\noalign{}
\begin{minipage}[b]{\linewidth}\raggedright
Nível
\end{minipage} & \begin{minipage}[b]{\linewidth}\raggedright
Exemplo
\end{minipage} & \begin{minipage}[b]{\linewidth}\raggedright
Detectável?
\end{minipage} & \begin{minipage}[b]{\linewidth}\raggedright
Relevância Ética
\end{minipage} \\
\midrule\noalign{}
\endhead
\bottomrule\noalign{}
\endlastfoot
1. Assistência linguística & Corretor ortográfico, paráfrase pontual &
Não & Mínima \\
2. Sumarização, tradução & DeepL, Elicit & Parcialmente & Baixa \\
3. Geração de seções & Redação de métodos ou discussão & Duvidosa &
Alta \\
4. Geração integral & Artigo completo por IA & Detectável (marcadores
grosseiros) & Gravíssima \\
\end{longtable}
}

A Portaria trata todos os níveis como equivalentes. Esta falta de
gradação expõe o pesquisador que usa Grammarly para correção ortográfica
ao mesmo regime legal de quem submete artigo integralmente gerado por IA
--- um tratamento juridicamente desproporcional. {[}\^{}8{]}: KOCAK, B.
et al.~Ensuring peer review integrity in the era of large language
models: a critical stocktaking. \emph{European Journal of Radiology
Artificial Intelligence}, v. 2, 100018, 2025. DOI:
10.1016/j.ejrai.2025.100018. Artigo de posicionamento com 47 co-autores
de 12 países. A taxonomia de níveis de uso de IA proposta pelos autores
(Seção 2) é referência internacional. Ponto forte: abordagem
multidisciplinar. Limitação: não oferece métricas de detecção por nível.
Relevância direta: a falta de gradação no art. 9º, alínea ``c'', ignora
esta distinção.

\subsubsection{\texorpdfstring{2.3. A falácia do
\emph{enforcement}}{2.3. A falácia do enforcement}}\label{a-faluxe1cia-do-enforcement}

A impossibilidade técnica de detecção produz uma \textbf{falácia de
enforcement}: o protocolo parece exigir conduta, mas não oferece meio de
verificação. Pérez-Neri et al.~(2025), na \emph{Frontiers in Artificial
Intelligence}, sintetizam: ``os desenvolvimentos atuais {[}em
detectores{]} não são suficientemente precisos ou confiáveis para uso em
avaliação acadêmica''{[}\^{}9{]}. {[}\^{}9{]}: PÉREZ-NERI, I. et
al.~Detecting and correcting errors in scientific literature in the
generative AI era. \emph{Frontiers in Artificial Intelligence}, v. 8,
2025. DOI: 10.3389/frai.2025.1644098. Revisão de escopo com análise de
22 ferramentas de detecção. Conclusão consistente com Weber-Wulff et
al.~(2023). Limitação: revisão narrativa, não sistemática. Forte: inclui
análise de ferramentas emergentes (2024-2025). Citação específica sobre
``insuficientemente precisos ou confiáveis'' está na Seção 3.2.

\subsection{3. Dimensão Hermenêutica: As Sete Ambiguidades do Art.
9º}\label{dimensuxe3o-hermenuxeautica-as-sete-ambiguidades-do-art.-9uxba}

O art. 9º da Portaria nº 2.664/2026 contém \textbf{pelo menos sete
ambiguidades} que comprometem sua aplicação coerente e produzem
insegurança jurídica: \textbar{} \# \textbar{} Ambiguidade \textbar{}
Dispositivo \textbar{} Problema Interpretativo \textbar{}
\textbar---\textbar-------------\textbar-------------\textbar------------------------\textbar{}
\textbar{} 1 \textbar{} ``pesquisa apoiada pelo CNPq'' \textbar{}
\emph{caput} \textbar{} Apenas financiada diretamente, ou qualquer
pesquisa de bolsista PQ? \textbar{} \textbar{} 2 \textbar{} ``Qualquer
espécie'' de IAG \textbar{} alínea ``c'' \textbar{} Inclui corretor
ortográfico? Tradutor automático? \textbar{} \textbar{} 3 \textbar{}
``Qualquer fase'' \textbar{} alínea ``c'' \textbar{} Concepção da ideia
já dispara a obrigação? \textbar{} \textbar{} 4 \textbar{} ``Conteúdo
gerado por IAG'' \textbar{} alínea ``d'' \textbar{} O que é ``gerado''?
E texto híbrido humano+máquina? \textbar{} \textbar{} 5 \textbar{}
``Finalidade'' a especificar \textbar{} alínea ``c'' \textbar{} Basta
``revisão textual'' ou exige detalhamento por parágrafo? \textbar{}
\textbar{} 6 \textbar{} ``Ferramenta utilizada'' \textbar{} alínea ``c''
\textbar{} Versão exata? Provedor? Data da consulta? \textbar{}
\textbar{} 7 \textbar{} Vedação a pareceres \textbar{} alínea ``e''
\textbar{} Aplica-se a LLM local (open-source) sem transmissão de dados?
\textbar{} A \textbf{ambiguidade \#1} é particularmente grave. A
expressão ``pesquisa apoiada pelo CNPq'' pode ser interpretada
restritivamente (apenas pesquisas com repasse financeiro direto) ou
ampliativamente (toda atividade de pesquisador que recebe bolsa PQ). A
primeira interpretação exclui parcela significativa da produção
científica brasileira; a segunda amplia o escopo além da competência
legal do CNPq, invadindo seara de periódicos e
universidades{[}\^{}10{]}.

A \textbf{ambiguidade \#4} --- o que constitui ``conteúdo gerado por
IAG'' --- é a mais problemática. Se o pesquisador usa ChatGPT para
sugerir três versões de um parágrafo, edita manualmente as três e funde
em uma quarta versão, o texto resultante é ``gerado por IAG''? A
literatura sobre \emph{human-AI co-creation} demonstra que a distinção é
um continuum, não uma dicotomia.

A \textbf{ambiguidade \#7} cria brecha facilmente explorável: IAGs
executadas localmente (Llama, DeepSeek, Qwen, Mistral) não transmitem
dados a terceiros. O CNPq não tem meios de auditar o ambiente
computacional de cada pesquisador-parecerista. A vedação, portanto, é
inócuas para usuários de modelos locais, que representam parcela
crescente da comunidade técnica\footnote{A ascensão de modelos
  open-source (Llama 3.1 405B, DeepSeek-V3, Mistral Large) viabiliza
  execução local com desempenho comparável a modelos comerciais. A
  alínea ``e'' não diferencia LLMs locais de APIs comerciais, criando
  incongruência: o mesmo uso de IA é vedado se acessado via ChatGPT
  (API), mas permitido virtualmente se executado localmente ---
  distinção sem relevância ética.}. {[}\^{}10{]}: Esta ambiguidade foi
identificada pelo Portal PesquisAI (2026) em análise publicada em
16/03/2026. O artigo questiona: ``tais diretrizes devem ser observadas
apenas em pesquisas financiadas pelo CNPq ou podem funcionar como
parâmetro normativo mais amplo?'' A ausência de resposta na Portaria
cria insegurança. PESQUISAI. CNPq reconhece o uso da Inteligência
Artificial Generativa na pesquisa científica: notas sobre a Portaria nº
2.664/2026. \emph{Portal PesquisAI}, 16 mar. 2026. Análise independente
com valor dogmático limitado por não ser publicação revisada por pares.

\subsection{4. Dimensão Ética: O Paradoxo da Transparência e o Dilema do
Inovador}\label{dimensuxe3o-uxe9tica-o-paradoxo-da-transparuxeancia-e-o-dilema-do-inovador}

\subsubsection{4.1. A transparência como
armadilha}\label{a-transparuxeancia-como-armadilha}

Schilke e Reimann (2025) identificaram o que denominam ``dilema da
transparência'' (\emph{transparency dilemma}) em estudo no
\emph{Organizational Behavior and Human Decision Processes}: a
declaração de uso de IA reduz a confiança do leitor no trabalho, mesmo
quando a IA foi usada de forma ética e assistiva{[}\^{}12{]}.

Este efeito é particularmente perverso no contexto do protocolo CNPq: o
pesquisador que declara uso de IA segundo o art. 9º, alínea ``c'',
expõe-se a escrutínio mais rigoroso, vieses de revisão e potencial
desqualificação --- mesmo que o uso tenha sido legítimo e assistivo. O
pesquisador que não declara, por outro lado, não corre risco (a detecção
é inviável). O resultado é a \textbf{seleção adversa}: o protocolo pune
a honestidade e recompensa a omissão. {[}\^{}12{]}: SCHILKE, O.;
REIMANN, M. The transparency dilemma: How AI disclosure erodes trust.
\emph{Organizational Behavior and Human Decision Processes}, v. 188,
2025. DOI: 10.1016/j.obhdp.2025.104405. Estudo experimental com 6
estudos (N total = 2.847). Revela que disclosure de IA reduz trust mesmo
quando a IA é usada apenas para formatação. Limitação: cenários
hipotéticos, não campo real. Relevância: sugere que a alínea ``c'' pode
ter efeito oposto ao desejado.

\subsubsection{4.2. O paradoxo da reprodutibilidade vs.~segredo
industrial}\label{o-paradoxo-da-reprodutibilidade-vs.-segredo-industrial}

A alínea ``d'' --- vedação de submissão de conteúdo gerado por IAG como
autoria humana --- colide com a realidade da pesquisa de fronteira
medicada por IA. Modelos de AlphaFold, GNoME ou descoberta por IA operam
sob lógica de caixa-preta protegida por segredo industrial.

Este paradoxo é bem descrito por Kocak et al.~(2025): a
reprodutibilidade científica exige transparência total, mas os modelos
algorítmicos que viabilizam inovações aceleradas são protegidos como
propriedade intelectual. A comunidade pode validar resultados finais,
não o processo cognitivo da máquina. A Portaria, ao tratar a
transparência como sempre possível, ignora o conflito estrutural entre
propriedade corporativa e escrutínio acadêmico{[}\^{}13{]}.
{[}\^{}13{]}: A tensão entre segredo industrial e integridade acadêmica
não é resolvida pela Portaria. O documento da CAPES (Brasil, 2026)
menciona o tema tangencialmente mas não oferece solução. Referência
complementar: MESZAROS, J.; HUYS, I.; IOANNIDIS, J. P. A. Challenges in
applying the EU AI Act research exemptions to contemporary AI research.
\emph{npj Digital Medicine}, 2026. DOI: 10.1038/s41746-025-02263-0.
Demonstra que o problema é global: a distinção entre ``pesquisa'' e
``atividade comercial'' no AI Act europeu é igualmente problemática.

\subsection{5. Dimensão Jurídico-Brasileira: Colisões
Normativas}\label{dimensuxe3o-juruxeddico-brasileira-colisuxf5es-normativas}

\subsubsection{5.1. LGPD vs.~alínea ``c''}\label{lgpd-vs.-aluxednea-c}

A alínea ``c'' exige que o pesquisador especifique ``a ferramenta
utilizada e a finalidade'' do uso de IAG. Se a ferramenta for um LLM
comercial (ChatGPT, Claude, Gemini), o pesquisador pode estar --- sem
saber --- violando a LGPD ao submeter dados de pesquisa a servidores no
exterior, especialmente se os dados envolverem informações pessoais de
sujeitos de pesquisa{[}\^{}14{]}.

A Portaria oferece zero orientação sobre como conciliar a obrigação de
declarar detalhadamente o uso de IA com a proteção de dados. O
pesquisador fica diante de um conflito normativo: cumprir o CNPq
(declarar a ferramenta) ou cumprir a LGPD (não expor dados a terceiros
não autorizados pelo consentimento). {[}\^{}14{]}: A LGPD (Lei nº
13.709/2018) exige consentimento específico e finalidade determinada
para tratamento de dados pessoais. A submissão de dados a LLMs
comerciais (que tipicamente usam dados para treinamento) pode configurar
violação. Cf. art. 7º, I c/c art. 11, I da LGPD. Adicionalmente, a
decisão do STF no RE 1.057.258 (Tema 987) reforça a autodeterminação
informativa. A Portaria do CNPq não menciona a LGPD --- lacuna grave.

\subsubsection{5.2. Art. 218 da Constituição
Federal}\label{art.-218-da-constituiuxe7uxe3o-federal}

O art. 218 da CF determina que ``o Estado promoverá e incentivará o
desenvolvimento científico, a pesquisa, a capacitação científica e
tecnológica e a inovação''. A Portaria impõe obrigações declaratórias e
vedações sem contrapartida de incentivo.

Este descompasso entre regulação e fomento pode ser questionado à luz do
princípio constitucional. A doutrina do direito administrativo da
inovação (Bucci, 2021) sustenta que a regulação estatal da pesquisa deve
observar a \textbf{proporcionalidade} entre meios restritivos e fins
promocionais{[}\^{}15{]}. Uma norma que cria custo burocrático sem
benefício demonstrável de integridade é desproporcional. {[}\^{}15{]}:
BUCCI, M. P. D. \emph{Fundamentos para uma teoria jurídica das políticas
públicas}. 2. ed.~São Paulo: Saraiva Jur, 2021. ISBN 9786555595802.
Defende que a regulação da pesquisa científica deve observar o princípio
da proporcionalidade na modalidade ``proibição de excesso''. A aplicação
ao caso concreto sugere que custos burocráticos sem contrapartida de
enforcement violam este princípio. Limitação: obra geral, não específica
sobre IA.

\subsubsection{5.3. Ausência de sanção específica e legalidade
estrita}\label{ausuxeancia-de-sanuxe7uxe3o-especuxedfica-e-legalidade-estrita}

A Portaria não estabelece sanções para o descumprimento do art. 9º. Esta
lacuna viola o princípio da legalidade estrita (art. 5º, II, CF) e a
tipicidade no direito administrativo sancionador. Sem descrição clara de
infração e sanção, a aplicação da norma dependeria de interpretação
extensiva --- vedada neste âmbito.

O direito administrativo brasileiro exige que a sanção esteja prevista
em lei ou ato normativo equivalente (art. 5º, XXXIX, CF c/c art. 1º da
LINDB). A ausência de sanção explícita significa que, na prática, a
Portaria é \textbf{lex imperfecta}: obriga mas não pune, ou pune sem
tipificação, o que é juridicamente insustentável{[}\^{}16{]}.
{[}\^{}16{]}: A doutrina de direito administrativo sancionador é clara:
não há sanção sem prévia cominação legal (nulla poena sine lege). A
Portaria CNPq é ato administrativo normativo, não lei em sentido formal.
Mesmo que se admita poder disciplinar da autarquia (art. 37, §6º, CF), a
ausência de especificação do que constitui infração e qual a penalidade
viola o devido processo legal substantivo. Ver: JUSTEN FILHO, M.
\emph{Curso de direito administrativo}. 15. ed.~Rio de Janeiro: Forense,
2024. ISBN 9786559649815.

\subsection{6. Dimensão Comparada: Quatro Modelos
Regulatórios}\label{dimensuxe3o-comparada-quatro-modelos-regulatuxf3rios}

A comparação internacional revela que cada país abordou o desafio de
forma diversa, e que a opção brasileira --- autodeclaração sem
verificação técnica --- é a de menor densidade normativa e menor
capacidade de enforcement.

\subsubsection{6.1. China: Regulação Centralizada com Padrão
Obrigatório}\label{china-regulauxe7uxe3o-centralizada-com-padruxe3o-obrigatuxf3rio}

A China adotou o modelo mais abrangente e tecnicamente sofisticado: -
\textbf{Ministério da Ciência e Tecnologia (2023)}: diretrizes proíbem
uso de IA generativa para gerar materiais de submissão, exigem rotulagem
clara de conteúdo gerado por IA e vedam IA como co-autora{[}\^{}17{]}. -
\textbf{Academia Chinesa de Ciências (2024)}: oito ``lembretes de
integridade'' detalhando usos permitidos e proibidos por etapa de
pesquisa. - \textbf{GB 45438-2025 (setembro de 2025)}: \textbf{padrão
nacional obrigatório} de marcação de conteúdo gerado por IA. Exige
marcação explícita (visível ao usuário) e implícita (metadados,
watermarks digitais)\footnote{CHINA. \emph{Measures for the
  Identification of AI-Generated (Synthetic) Content}. Promulgado
  conjuntamente por CAC, MIIT, MPS e NRTA em 07/03/2025, vigente desde
  01/09/2025. Acompanhado do padrão nacional GB 45438-2025
  (\emph{Cybersecurity Technology: Labeling Method for Content Generated
  by Artificial Intelligence}). Referência:
  https://regulations.ai/regulations/RAI-CN-NA-MIASCXX-2025. Ponto
  forte: estabelece padrão técnico objetivo, não autodeclaração.
  Fraqueza: aplica-se apenas a provedores de IA chineses, não a usuários
  finais.}. - \textbf{Guideline on Boundaries of AIGC Usage 3.0
(dezembro de 2025)}: 3ª edição com casos práticos, especificando limites
por etapa (revisão de literatura, coleta de dados, análise estatística,
geração de figuras).

O modelo chinês é o que mais se aproxima de um \emph{enforcement}
factível, por três razões: (a) o padrão técnico obrigatório (GB
45438-2025) é aplicado na origem pelos provedores de IA chineses
(aplicam marcação automática); (b) a fiscalização é exercida pelo CAC
(Cyberspace Administration) com poder de sanção administrativa; (c) o
regime exige registro e filing de algoritmos de IA.

Limitações do modelo chinês: (a) dependência de provedores nacionais ---
LLMs estrangeiros (ChatGPT, Claude) não são acessíveis na China
continental, o que cria ecossistema controlado; (b) alto custo de
conformidade; (c) potencial para censura e controle político disfarçado
de regulação de IA\footnote{A centralização regulatória chinesa tem
  dupla face: por um lado, permite coordenação e escala; por outro,
  subordina a integridade acadêmica aos interesses geopolíticos do
  Partido-Estado. A exigência de ``valores socialistas centrais'' na
  geração de conteúdo (art. 4º dos \emph{Interim Measures on Generative
  AI Services}) cria filtro de conteúdo que a academia chinesa
  internaliza. Ver: BELADIYA, R. How China Is Controlling Generative AI
  Technologies. \emph{AI Law Guide}, maio 2026.}. {[}\^{}17{]}: CHINA.
Ministry of Science and Technology. \emph{Guidelines for Responsible
Research Conduct}. Beijing, 2023. Publicado em 21/12/2023. Texto
integral disponível em:
http://english.www.gov.cn/news/202401/06/content\_WS6598c927c6d0868f4e8e2d24.html.
Acesso em: 15 maio 2026. Inovação: exige que conteúdo gerado por IA seja
claramente rotulado e que o processo de geração seja explicado.
Limitação: diretrizes, não lei; enforcement depende de incorporação por
instituições.

\subsubsection{6.2. União Europeia: Abordagem Baseada em Risco com
Enforcement
Escalonado}\label{uniuxe3o-europeia-abordagem-baseada-em-risco-com-enforcement-escalonado}

O EU AI Act (Regulation 2024/1689) é o primeiro marco regulatório
abrangente de IA no mundo. Para a pesquisa acadêmica, seus dispositivos
mais relevantes são: - \textbf{Art. 2(6) e (8)}: isenções para pesquisa
(``research exemption'') --- sistemas de IA desenvolvidos exclusivamente
para pesquisa científica estão fora do escopo{[}\^{}20{]}. -
\textbf{Art. 50}: obrigações de transparência aplicáveis a todos os
sistemas de IA (incluindo pesquisa aplicada) --- exigem que conteúdo
sintético seja marcado e detectável\footnote{Art. 50(2): ``Providers of
  GPAI models that generate synthetic content shall ensure that the
  output is marked in a machine-readable format and detectable as
  artificially generated or manipulated.'' O art. 50(4) adiciona
  requisitos específicos para \emph{deepfakes} e para texto de interesse
  público. Estas obrigações entram em vigor em 2 de agosto de 2026.
  GILS, T. et al.~From Policy to Practice: Prototyping The EU AI Act's
  Transparency Requirements. \emph{SSRN}, 2024. DOI:
  10.2139/ssrn.4714345. Estudo de policy prototyping que testou os
  requisitos na prática e identificou lacunas na implementação.}. -
\textbf{Code of Practice on Marking and Labelling} (2º draft em
05/03/2026, final esperado para junho de 2026): especificações técnicas
para marcação de conteúdo gerado por IA, aplicáveis a partir de agosto
de 2026. - \textbf{Diretrizes da Comissão Europeia} (2º trimestre de
2026): interpretação oficial das obrigações de transparência.

A distinção crítica do modelo europeu é: não impõe declaração pelo autor
--- impõe \textbf{marcação técnica pelo provedor}. O LLM comercial é
obrigado por lei (art. 50(2)) a marcar seu output de forma legível por
máquina. A obrigação de transparência recai sobre quem
\emph{disponibiliza} o sistema, não sobre quem o \emph{usa}.

Isso contrasta fundamentalmente com o modelo brasileiro, que impõe a
obrigação exclusivamente sobre o pesquisador-usuário --- sem exigência
correspondente sobre os provedores. A assimetria é evidente: o Brasil
regula quem não tem como verificar a própria declaração.

Meszaros, Huys e Ioannidis (2026) identificaram fragilidades nas
isenções de pesquisa do AI Act: a distinção entre pesquisa acadêmica e
atividade comercial é turva; o desenvolvimento em ambiente controlado
(\emph{sandbox}) versus teste em ambiente real é difícil de
determinar\footnote{MESZAROS, J.; HUYS, I.; IOANNIDIS, J. P. A.
  Challenges in applying the EU AI Act research exemptions to
  contemporary AI research. \emph{npj Digital Medicine}, 2026. DOI:
  10.1038/s41746-025-02263-0. Análise jurídico-técnica rigorosa. Conclui
  que as isenções de pesquisa do AI Act ``rely on distinctions that may
  not fully capture the realities of contemporary AI research''.
  Relevância: mostra que mesmo o marco mais sofisticado do mundo tem
  limitações. Limitação: análise focada em pesquisa médica, pode não
  generalizar.}. {[}\^{}20{]}: EUROPEAN UNION. Regulation (EU) 2024/1689
of the European Parliament and of the Council of 13 June 2024 laying
down harmonised rules on artificial intelligence (Artificial
Intelligence Act). \emph{Official Journal of the European Union},
Luxembourg, 2024. Art. 2(6): ``This Regulation does not apply to AI
systems specifically developed and put into service for the sole purpose
of scientific research and development.'' A isenção é generosa mas
limitada a sistemas desenvolvidos \emph{exclusivamente} para pesquisa
--- não cobre uso de LLMs comerciais para redação acadêmica. Texto
completo: https://eur-lex.europa.eu/eli/reg/2024/1689.

\subsubsection{6.3. Estados Unidos: Segurança de pesquisa em vez de
Regulação de
IA}\label{estados-unidos-seguranuxe7a-de-pesquisa-em-vez-de-regulauxe7uxe3o-de-ia}

Os EUA não possuem regulação federal específica sobre uso de IA na
pesquisa acadêmica. A abordagem é indireta, centrada em segurança de
pesquisa (\emph{research security}) e combate à influência estrangeira:
- \textbf{NSF Scientific Integrity Policy (NSF 24-007, fevereiro de
2024)}: reafirma compromisso com integridade, mas não menciona IA
especificamente{[}\^{}23{]}. - \textbf{NSPM-33 e TRUST Framework}: foco
em \emph{disclosure} de apoio estrangeiro e conflitos de compromisso,
não em declaração de IA\footnote{NSF. \emph{Trusted Research Using
  Safeguards and Transparency (TRUST)}. Jun.~2024. Disponível em:
  https://nsf-gov-resources.nsf.gov/files/NSF\%20OCRSSP\%20TRUST\%20Policy\%20Memo.pdf.
  Framework de árvore decisória para avaliar riscos de segurança
  nacional. Três branches: (1) appointments ativos e suporte à pesquisa,
  (2) não-disclosure, (3) considerações de segurança nacional. Nenhum
  branch trata de IA. Foco: influência estrangeira, não integridade da
  escrita acadêmica.}. - \textbf{CHIPS and Science Act (2022)}: §10339B
exige \emph{disclosure} de apoio financeiro estrangeiro, mas silencia
sobre IA. - \textbf{NSF Proposal \& Award Policies (PAPPG 24-1, maio de
2024)}: incorpora requisitos de segurança de pesquisa e integridade
científica, mas IA não é tratada como categoria específica\footnote{NSF.
  \emph{Proposal and Award Policies and Procedures Guide} (PAPPG 24-1).
  Maio 2024. Cap. XI.M, Scientific Integrity. Disponível em:
  https://new.nsf.gov/policies/pappg. Incorpora definição do NSTC. A
  PAPPG foi suplementada por policy notice em dez. 2025 incorporando
  §10339B (CHIPS Act) sobre disclosure de apoio financeiro estrangeiro.
  Em nenhum ponto há menção a IA generativa. Isto significa zero custo
  de conformidade para pesquisadores americanos --- e zero proteção
  contra má-uso.}.

O modelo americano é o mais liberal: não exige declaração de IA, não
estabelece padrão de marcação, não cria obrigação para provedores. A
integridade é delegada às instituições e periódicos, que têm políticas
próprias. Esta abordagem minimiza custos burocráticos mas transfere o
problema para o \emph{self-governance} de cada entidade. {[}\^{}23{]}:
NSF. \emph{NSF Scientific Integrity Policy} (NSF 24-007). 12 fev. 2024.
Disponível em:
https://nsf-gov-resources.nsf.gov/pubs/2024/nsf24007/nsf24007.pdf.
Política genérica que ``estabelece expectativas para manter e aprimorar
a Integridade Científica na NSF''. Silêncio sobre IA reflete a posição
do governo Biden: a ordem executiva de outubro de 2023 (EO 14110) sobre
IA segura e confiável focou em riscos sistêmicos, não em escrita
acadêmica. A gestão Trump (2025-) pode alterar esta abordagem.

\subsubsection{6.4. Tabela Comparativa}\label{tabela-comparativa}

{\def\LTcaptype{none} % do not increment counter
\begin{longtable}[]{@{}
  >{\raggedright\arraybackslash}p{(\linewidth - 8\tabcolsep) * \real{0.1429}}
  >{\raggedright\arraybackslash}p{(\linewidth - 8\tabcolsep) * \real{0.3714}}
  >{\raggedright\arraybackslash}p{(\linewidth - 8\tabcolsep) * \real{0.1000}}
  >{\raggedright\arraybackslash}p{(\linewidth - 8\tabcolsep) * \real{0.2286}}
  >{\raggedright\arraybackslash}p{(\linewidth - 8\tabcolsep) * \real{0.1571}}@{}}
\toprule\noalign{}
\begin{minipage}[b]{\linewidth}\raggedright
Dimensão
\end{minipage} & \begin{minipage}[b]{\linewidth}\raggedright
Brasil (CNPq 2.664/2026)
\end{minipage} & \begin{minipage}[b]{\linewidth}\raggedright
China
\end{minipage} & \begin{minipage}[b]{\linewidth}\raggedright
União Europeia
\end{minipage} & \begin{minipage}[b]{\linewidth}\raggedright
EUA (NSF)
\end{minipage} \\
\midrule\noalign{}
\endhead
\bottomrule\noalign{}
\endlastfoot
Base legal & Portaria administrativa & Lei + Padrão Nacional &
Regulamento (diretamente aplicável) & Política interna NSF \\
Obrigação & Autodeclaração do pesquisador & Marcação pelo provedor +
declaração do usuário & Marcação técnica pelo provedor (Art. 50) &
Nenhuma específica \\
Padrão técnico & Nenhum & GB 45438-2025 (obrigatório) & Code of Practice
(voluntário) & Nenhum \\
Enforcement & Inexistente (sem sanção) & CAC com sanções administrativas
& Comissão + autoridades nacionais & pesquisa segurança (não IA) \\
Escopo & pesquisa apoiada pelo CNPq & Toda pesquisa na China & Sistemas
de IA no mercado EU & Pesquisas financiadas NSF \\
Custo de conformidade para pesquisador & Alto (declaração detalhada) &
Médio (rotulagem + verificação) & Baixo (marcação é do provedor) &
Zero \\
\end{longtable}
}

\subsubsection{6.5. Síntese Comparativa}\label{suxedntese-comparativa}

Cada modelo reflete escolhas políticas e históricas distintas. O Brasil
optou pelo \textbf{modelo de autodeclaração com custo concentrado no
pesquisador} --- a pior combinação possível: alto custo de conformidade
para quem cumpre, zero custo para quem descumpre, e impossibilidade
técnica de verificação.

A China optou pelo \textbf{modelo de padrão técnico obrigatório na
fonte} --- mais eficaz, mas dependente de controle centralizado sobre
provedores de IA e compatível com regime de vigilância digital.

A UE optou pelo \textbf{modelo de regulação híbrida} --- impõe
obrigações sobre provedores (marcação) e sobre usuários (transparência
em certos usos), com abordagem baseada em risco.

Os EUA optaram pelo \textbf{modelo de não-regulação específica} ---
trata IA como ferramenta de pesquisa, não como categoria regulatória
distinta, e foca segurança em vez de integridade textual.

A Portaria CNPq nº 2.664/2026, neste cenário, é o único instrumento que
impõe obrigação sem oferecer meio de cumprimento. É regulação sem
tecnologia.

\subsection{7. Dimensão Estrutural: O Sistema Qualis-CAPES como
Antecedente}\label{dimensuxe3o-estrutural-o-sistema-qualis-capes-como-antecedente}

A Portaria nº 2.664/2026 não surge em vácuo. Ela se insere em tradição
de avaliação quantitativa e burocrática que sistematicamente privilegia
métricas internacionais em detrimento da pesquisa nacional --- como
demonstrado pela Nota Técnica da AGB (2024){[}\^{}26{]} e pelo estudo de
Caballero Rivero et al.~(2024)\footnote{CABALLERO RIVERO, A.; SANTOS, R.
  N. M.; TRZESNIAK, P. Efeitos dos sistemas de avaliação de pesquisa de
  CAPES e CNPQ nos padrões de publicação dos pesquisadores das ciências
  da saúde no Brasil. \emph{Em Questão}, Porto Alegre, v. 30, 2024. DOI:
  10.1590/1808-5245.30.138437. Qualis A1. Metodologia: questionário
  semiestruturado com amostra estratificada. Demonstra que critérios
  avaliativos distorcem escolhas de publicação. Relevância direta para
  este ensaio: o mesmo efeito performativo se aplicará à declaração de
  IA.}.

O efeito performativo é bem descrito por Albuquerque (2025): ``o simples
fato de existirem critérios externos faz com que os programas de
pós-graduação se reorganizem para atendê-los''\footnote{ALBUQUERQUE, U.
  P. Ciência em debate: o peso silencioso da avaliação de cursos e
  publicações científicas. \emph{The Conversation}, 13 out. 2025. Texto
  de divulgação de pesquisador sênior (UFPE). Embora não seja artigo
  revisado, sintetiza pesquisa acumulada do autor sobre o tema. Citação
  direta sobre ``efeito performativo'' no parágrafo 7.}. Sguissardi
(2006) denominou este fenômeno de ``avaliação defensiva''.

Perez (2025) atualizou a crítica para o contexto do Qualis 2025-2028,
argumentando que os novos critérios da CAPES ``tendem a aprofundar
desigualdades regionais, institucionais e de gênero''\footnote{PEREZ, O.
  C. Avaliação em Disputa: A Reforma do Qualis e os Desafios para a
  Ciência. \emph{Novos Debates}, v. 11, n.~1, 2025. DOI:
  10.48006/2358-0097/V11N1.E111011. Qualis A2. Limitação: ensaio teórico
  sem dados primários. Contribuição: aponta que as novas regras
  ``penalizam programas novos e pesquisadores em início de carreira''.}.

O custo operacional do novo modelo multidimensional é estimado em R\$
150 milhões iniciais + R\$ 50 milhões/ano (CAPES, 2024)\footnote{CAPES.
  Ofício Circular nº 46/2024: diretrizes para o ciclo avaliativo
  2025-2028. Brasília, 2024. A estimativa de R\$ 150M é citada em
  análise crítica da LinkedIn (Costa, 2024) e em reportagem da ADUFRJ
  (11/12/2024). Dados oficiais consolidados não foram publicados.}.
{[}\^{}26{]}: AGB. Nota técnica: reflexões sobre o Qualis CAPES, o fator
de impacto e a divulgação científica. São Paulo, 2024. Disponível em:
https://agb.org.br/nota-tecnica-reflexoes-sobre-o-qualis-capes-o-fator-de-impacto-e-a-divulgacao-cientifica/.
Crítica contundente ao viés estrangeiro do Qualis. Limitação: documento
de associação, não artigo revisado por pares.

\subsection{8. Balanço Crítico: Acertos, Fragilidades e
Recomendações}\label{balanuxe7o-cruxedtico-acertos-fragilidades-e-recomendauxe7uxf5es}

\subsubsection{8.1. Acertos da Portaria}\label{acertos-da-portaria}

\begin{enumerate}
\def\labelenumi{\arabic{enumi}.}
\tightlist
\item
  \textbf{Reconhecimento institucional}: a Portaria tira a IAG da
  invisibilidade normativa e a reconhece como fenômeno relevante na
  pesquisa --- mérito não trivial.
\item
  \textbf{Não proibição}: diferentemente de posições reacionárias, a
  Portaria não proíbe IA; busca regular uso responsável.
\item
  \textbf{Atualização do debate}: acelera a discussão sobre integridade
  em universidades e periódicos.
\item
  \textbf{Alinhamento internacional}: ainda que imperfeita, a Portaria
  se insere em movimento global de regulação.
\end{enumerate}

\subsubsection{8.2. Fragilidades}\label{fragilidades}

\begin{enumerate}
\def\labelenumi{\arabic{enumi}.}
\tightlist
\item
  \textbf{Inaplicabilidade técnica}: detectores são falíveis; a
  proporção 40:1 invalida a premissa de enforcement.
\item
  \textbf{Custo sem benefício}: impõe burocracia ao pesquisador sem
  oferecer ganho demonstrável de integridade.
\item
  \textbf{Assimetria punitiva}: penaliza quem declara, não quem omite.
\item
  \textbf{Ambiguidade normativa}: 7 ambiguidades no art. 9º geram
  insegurança jurídica.
\item
  \textbf{Colisão com LGPD e CF}: conflitos normativos não resolvidos.
\item
  \textbf{Inexequibilidade da alínea ``e''}: vedação a pareceres com IAG
  é facilmente contornada por modelos locais.
\item
  \textbf{Sem padrão técnico}: diferentemente de China (GB 45438-2025) e
  UE (Art. 50 AI Act), o Brasil não exige marcação técnica pelos
  provedores.
\end{enumerate}

\subsubsection{8.3. Recomendações}\label{recomendauxe7uxf5es}

\begin{enumerate}
\def\labelenumi{\arabic{enumi}.}
\tightlist
\item
  \textbf{Adotar modelo híbrido} nos moldes da UE: obrigação de marcação
  técnica pelos provedores de IA (LLMs comerciais) somada à declaração
  simplificada pelo pesquisador.
\item
  \textbf{Graduar a exigência} por nível de risco (assistência
  linguística vs.~geração substantiva).
\item
  \textbf{Substituir declaração detalhada} por \emph{checklist}
  padronizado de uma página.
\item
  \textbf{Estabelecer sanções claras} respeitando o princípio da
  legalidade.
\item
  \textbf{Harmonizar com a LGPD} e com o art. 218 CF.
\item
  \textbf{Negociar com provedores de IA} a marcação automática de
  conteúdo gerado em português.
\end{enumerate}

\subsection{9. Conclusão: A Regulação que a Tecnologia Ainda Não
Alcançou}\label{conclusuxe3o-a-regulauxe7uxe3o-que-a-tecnologia-ainda-nuxe3o-alcanuxe7ou}

Este ensaio demonstrou que a Portaria CNPq nº 2.664/2026, não obstante
suas boas intenções, opera como instrumento de retardamento burocrático
da pesquisa científica brasileira. As evidências são convergentes:
detectores falham, a proporção uso/declaração é de 40:1, o texto contém
ambiguidades insanáveis, a fiscalização é inviável, e o modelo
brasileiro é o que oferece menor capacidade de enforcement entre as
quatro principais jurisdições analisadas.

O paradoxo central é que a Portaria tenta resolver um problema técnico
--- a detecção de IA --- com instrumento jurídico --- a declaração
obrigatória. Este descompasso entre a natureza do problema e o meio de
solução condena o protocolo à ineficácia.

O pesquisador que declarar o uso de IA será escrutinado; o que não
declarar, não. O inovador honesto paga o custo burocrático; o
mal-intencionado paga zero. A Portaria não protege a integridade: apenas
a burocratiza.

\subsection{A saída não é abandonar a integridade, mas redefini-la.
Avaliação substantiva baseada em impacto replicável e relevância social
será sempre mais eficaz que declaratórias inverificáveis. Enquanto o
CNPq não enfrentar o problema real --- a inviabilidade técnica da
detecção ---, a ciência brasileira continuará pagando o preço de um
protocolo que não protege, mas apenas
retarda.}\label{a-sauxedda-nuxe3o-uxe9-abandonar-a-integridade-mas-redefini-la.-avaliauxe7uxe3o-substantiva-baseada-em-impacto-replicuxe1vel-e-relevuxe2ncia-social-seruxe1-sempre-mais-eficaz-que-declaratuxf3rias-inverificuxe1veis.-enquanto-o-cnpq-nuxe3o-enfrentar-o-problema-real-a-inviabilidade-tuxe9cnica-da-detecuxe7uxe3o-a-ciuxeancia-brasileira-continuaruxe1-pagando-o-preuxe7o-de-um-protocolo-que-nuxe3o-protege-mas-apenas-retarda.}

\section{Referências}\label{referuxeancias}

AGB --- ASSOCIAÇÃO DOS GEÓGRAFOS BRASILEIROS. Nota técnica: reflexões
sobre o Qualis CAPES, o fator de impacto e a divulgação científica. São
Paulo, 2024. Disponível em:
https://agb.org.br/nota-tecnica-reflexoes-sobre-o-qualis-capes-o-fator-de-impacto-e-a-divulgacao-cientifica/
ALBUQUERQUE, U. P. Ciência em debate: o peso silencioso da avaliação de
cursos e publicações científicas. \textbf{The Conversation}, 13 out.
2025. Disponível em:
https://theconversation.com/ciencia-em-debate-o-peso-silencioso-da-avaliação-de-cursos-e-publicacoes-cientificas-266875
BELADIYA, R. How China Is Controlling Generative AI Technologies.
\textbf{AI Law Guide}, maio 2026. Disponível em:
https://ailawguide.org/blog/how-china-is-controlling-generative-ai-technologies
BRASIL. Constituição da República Federativa do Brasil de 1988. Art. 5º,
II e XXXIX; Art. 218. BRASIL. Lei nº 13.709, de 14 de agosto de 2018
(Lei Geral de Proteção de Dados Pessoais). BRASIL. Portaria CNPq nº
2.664, de 6 de março de 2026. Institui a Política de Integridade na
Atividade Científica. \textbf{Diário Oficial da União}, 11 mar. 2026.
Disponível em:
http://memoria2.cnpq.br/web/guest/view/-/journal\_content/56\_INSTANCE\_0oED/10157/23142775
BRASIL, A. L. \textbf{A inteligência artificial na pesquisa e no
fomento}: desafios e oportunidades. Brasília: CAPES, 2026. BUCCI, M. P.
D. \textbf{Fundamentos para uma teoria jurídica das políticas públicas}.
2. ed.~São Paulo: Saraiva Jur, 2021. CABALLERO RIVERO, A.; SANTOS, R. N.
M.; TRZESNIAK, P. Efeitos dos sistemas de avaliação de pesquisa de CAPES
e CNPQ nos padrões de publicação dos pesquisadores das ciências da saúde
no Brasil. \textbf{Em Questão}, Porto Alegre, v. 30, 2024. DOI:
10.1590/1808-5245.30.138437 CAPES. Ofício Circular nº 46/2024.
Diretrizes para o ciclo avaliativo 2025-2028. Brasília, 2024. CHINA.
Ministry of Science and Technology. Guidelines for Responsible Research
Conduct. Beijing, 2023. Disponível em:
http://english.www.gov.cn/news/202401/06/content\_WS6598c927c6d0868f4e8e2d24.html
CHINA. Measures for the Identification of AI-Generated (Synthetic)
Content. CAC/MIIT/MPS/NRTA, 2025. Vigente desde 01/09/2025. Padrão GB
45438-2025. EUROPEAN UNION. Regulation (EU) 2024/1689 (Artificial
Intelligence Act). \textbf{Official Journal of the European Union},
2024. Disponível em: https://eur-lex.europa.eu/eli/reg/2024/1689 GILS,
T. et al.~From Policy to Practice: Prototyping The EU AI Act's
Transparency Requirements. \textbf{SSRN}, 2024. DOI:
10.2139/ssrn.4714345 KOCAK, B. et al.~Ensuring peer review integrity in
the era of large language models. \textbf{European Journal of Radiology
Artificial Intelligence}, v. 2, 100018, 2025. DOI:
10.1016/j.ejrai.2025.100018 KWON, D. Science sleuths flag hundreds of
papers that use AI without disclosing it. \textbf{Nature}, v. 641,
p.~290-291, 2025. DOI: 10.1038/d41586-025-01180-2 LIANG, W. et
al.~Discrepancy between AI content proportion and disclosure rate in
scientific publications. \textbf{Nature Human Behaviour}, 2025. DOI:
10.48550/arXiv.2512.06705 LIU, J. Q. J. et al.~The great detectives:
humans versus AI detectors in catching large language model-generated
medical writing. \textbf{International Journal for Educational
Integrity}, v. 20, n.~1, 2024. DOI: 10.1007/s40979-024-00155-6 MESZAROS,
J.; HUYS, I.; IOANNIDIS, J. P. A. Challenges in applying the EU AI Act
research exemptions to contemporary AI research. \textbf{npj Digital
Medicine}, 2026. DOI: 10.1038/s41746-025-02263-0 NSF. NSF Scientific
Integrity Policy (NSF 24-007). 2024. Disponível em:
https://nsf-gov-resources.nsf.gov/pubs/2024/nsf24007/nsf24007.pdf NSF.
Trusted Research Using Safeguards and Transparency (TRUST). 2024.
Disponível em:
https://nsf-gov-resources.nsf.gov/files/NSF\%20OCRSSP\%20TRUST\%20Policy\%20Memo.pdf
PEREZ, O. C. Avaliação em Disputa: A Reforma do Qualis e os Desafios
para a Ciência. \textbf{Novos Debates}, v. 11, n.~1, 2025. DOI:
10.48006/2358-0097/V11N1.E111011 PÉREZ-NERI, I. et al.~Detecting and
correcting errors in scientific literature in the generative AI era.
\textbf{Frontiers in Artificial Intelligence}, v. 8, 2025. DOI:
10.3389/frai.2025.1644098 PESQUISAI. CNPq reconhece o uso da
Inteligência Artificial Generativa na pesquisa científica.
\textbf{Portal PesquisAI}, 16 mar. 2026. Disponível em:
https://www.pesquisai.info/post/cnpq-reconhece-o-uso-da-intelig\%C3\%AAncia-artificial-generativa-na-pesquisa-cient\%C3\%ADfica-notas-sobre-a-por
PUDASAINI, S. et al.~Survey on AI-generated plagiarism detection.
\textbf{Journal of Academic Ethics}, v. 23, p.~1137-1170, 2025. DOI:
10.1007/s10805-024-09576-x SCHILKE, O.; REIMANN, M. The transparency
dilemma: How AI disclosure erodes trust. \textbf{Organizational Behavior
and Human Decision Processes}, v. 188, 2025. DOI:
10.1016/j.obhdp.2025.104405 SGUISSARDI, V. A avaliação defensiva no
``modelo CAPES de avaliação''. \textbf{Perspectiva}, Florianópolis, v.
24, n.~1, p.~49-88, 2006. Disponível em:
https://repositorio.ufsc.br/handle/123456789/186529 WEBER-WULFF, D. et
al.~Testing of detection tools for AI-generated text.
\textbf{International Journal for Educational Integrity}, v. 19, n.~1,
2023. DOI: 10.1007/s40979-023-00146-z

\end{document}
